\section{Source lock and conventions}

All Hilbert-space inner products are conjugate-linear in the first slot.  The
literal clustered coefficient inherited from the V59 source chain is
\[
C_h=\sum_{\substack{d\in\mathcal D_x\\h\mid d}}
       \frac{\mu(d)\log d}{d},
\qquad
H=x^{21/32},\quad Q=x^{1/3},\quad U=x^{133/400}.
\]
It is real.  Reduced rational-frequency clustering produces one coefficient
per primitive denominator block, and primitive representatives give canonical
exactly-once coefficient coordinates.  The frozen packet kernels retain the
same outer $C_h$ for every phase label.

Those facts do not by themselves produce two source-native lanes.  In
particular, the current source compiler leaves open a literal map from the V59
$\beta,w$ sequences to primitive vectors $b_h,w_h$ in a common synthesis
space.  We therefore separate the unconditional finite-dimensional theorems
from their conditional V59 interpretation.

For a finite index set $\cA$, write
\[
\cH=\bigoplus_{h\in\cA}\cH_h.
\]
The direct sum is orthogonal.  Later, $J_h:\cH_h\to\cK$ denotes an arbitrary
linear embedding into one ambient Hilbert space; it is not assumed isometric.

An ``outer phase change'' means $C_h\mapsto\eta_hC_h$ with
$|\eta_h|=1$ after $C_h$ has been formed.  It is not a change of any individual
Möbius summand inside $C_h$.
